% Propietario conservado: /Users/ruben/Documents/ChatGPT/jueces y controles/REVISION_ARGUMENTAL_APERTURAS_CIERRES_20260915/II/manuscript_es/sections/05_accion.tex
\section{Constantes de Planck y estructura determinantal de la acción}
\label{ii:sec:action}
\subsection{Escala determinantal de acción y valoración decimal}

La escala absoluta y el corrector relativo se originan en objetos
distintos. El primero utiliza el bloque local del registro dodecafásico
de paso tres,
\begin{equation}
 H_4=\begin{pmatrix}
 1&1&1&1\\1&1&-1&-1\\1&-1&1&-1\\1&-1&-1&1
 \end{pmatrix},\qquad
 G_H=\mathbb Z^4/H_4\mathbb Z^4.
 \label{ii:eq:el-h4}
\end{equation}
La operación utiliza un cociente local. Los tres bloques del registro
completo no se sustituyen por una tensorización que multiplicara de
nuevo el número de hojas de este lector.

\begin{lemma}[Número de hojas del cociente local]
La forma normal de Smith de \(H_4\) es
\(\operatorname{diag}(1,2,2,4)\). En consecuencia,
\[
 G_H\simeq\mathbb Z/2\mathbb Z\oplus\mathbb Z/2\mathbb Z
 \oplus\mathbb Z/4\mathbb Z,\qquad |G_H|=16.
\]
\end{lemma}
\begin{proof}
El máximo común divisor de las entradas es uno. Los menores de orden dos
son pares y hay uno de valor \(-2\), de modo que su máximo común divisor
es dos. Como \(H_4H_4^{\mathsf T}=4I\) y \(\det H_4=-16\), los
cofactores son \(\pm4\); el tercer divisor determinantal es cuatro.
El cuarto es \(16\). Los cocientes sucesivos de los divisores
\(1,2,4,16\) son \(1,2,2,4\). Su producto da el orden del cociente.
\end{proof}

\begin{definition}[Lector determinantal de acción]
La norma multiplicativa de la sección escalar \(\alpha\), constante
sobre las hojas de \(G_H\), y una aplicación de la autoescala
\(\varphi\) producen
\begin{equation}
 \Sigma_H=\varphi\,\mathrm N_{G_H}(\alpha),\qquad
 \mathrm N_{G_H}(\alpha)=\prod_{g\in G_H}\alpha=\alpha^{16}.
 \label{ii:eq:el-lector-determinantal}
\end{equation}
La carta decimal asociada determina
\(\operatorname{ord}_{10}(\Sigma_H)=\lfloor\log_{10}\Sigma_H\rfloor\).
\end{definition}

La potencia dieciséis procede así del número de hojas del cociente, una
vez fijado este lector local. El factor \(\varphi\) actúa una vez sobre
la base de la recta de acción. El cálculo de Smith, por sí solo, no
seleccionaría cualquier funcional posible sobre el cociente: la norma
multiplicativa y la autoescala forman parte de la operación declarada.

\begin{theorem}[Década de la sección de acción]
El lector \eqref{ii:eq:el-lector-determinantal}, aplicado a las coordenadas
HMT ya generadas, satisface
\begin{equation}
 10^{-34}<\Sigma_H<10^{-33},\qquad
 \operatorname{ord}_{10}(\Sigma_H)=-34.
 \label{ii:eq:el-decada}
\end{equation}
\end{theorem}
\begin{proof}
Bastan los intervalos racionales, mucho más gruesos que los cilindros
disponibles,
\[
 \frac{1618}{1000}<\varphi<\frac{1619}{1000},\qquad
 \frac{7297}{10^6}<\alpha<\frac{7298}{10^6}.
\]
La función \((x,a)\mapsto xa^{16}\) es creciente en ambos argumentos
positivos. Las dos desigualdades enteras
\[
 1618\cdot7297^{16}>10^{65},\qquad
 1619\cdot7298^{16}<10^{66}
\]
implican
\[
 10^{-34}
 <\frac{1618}{1000}\left(\frac{7297}{10^6}\right)^{16}
 <\Sigma_H
 <\frac{1619}{1000}\left(\frac{7298}{10^6}\right)^{16}
 <10^{-33}.
\]
La definición del orden decimal concluye la prueba. No interviene una
unidad SI ni una masa de referencia en estas desigualdades.
\end{proof}

La sección normalizada comienza por
\(\Sigma_H=1.0462438381\ldots\times10^{-34}\). Su mantisa no se
identifica con la coordenada de acción anterior o posterior al retorno.
El lector determinantal fija el orden decimal; las dos coordenadas de
retorno se obtienen mediante los lectores siguientes.


\subsection{Secciones de acción y cociente adimensional de retorno}

Sea \(\mathcal L_S\) una recta orientada de acción y \(\mathcal U_S\)
su base positiva. El carácter de quinto orden y el término regional de
retorno tienen las expresiones
\begin{align}
 H_5={}&\frac12\sqrt{\frac{1000\alpha}{\varphi}}-\alpha
   +\frac9{16}\alpha^2-\frac59\alpha^3
   +\frac7{48}\alpha^4-\frac1{54}\alpha^5,
   \label{ii:eq:el-h5}\\
 D_A={}&\exp\!\left(-\frac{100\pi\alpha}{9}\right),
   \label{ii:eq:el-da}\\
 \mathcal C_\pi(\alpha)={}&169\alpha^6+
       \frac{D_A\alpha^7}{1-D_A\alpha},
   \label{ii:eq:el-cpi}\\
 \eta_{\mathrm{ret}}={}&H_5-\frac{90}{\pi}\mathcal C_\pi(\alpha).
   \label{ii:eq:el-eta}
\end{align}
Estas son lecturas posteriores del cierre dodecafásico y de su
incidencia regional. El entero \(169\) se conserva como el selector
regional ya construido, y los coeficientes de \(H_5\) como los de su
carácter local; no se reemplazan por una interpolación de la constante de
Planck. Las fórmulas explicitan exactamente qué lecturas se necesitan
para la etapa electrónica. No es necesario introducir aquí otros
funcionales de la teoría de masas.

Para \(0<\alpha<1\), se tiene \(0<D_A<1\), de modo que
\(D_A\alpha<1\). El término racional de \eqref{ii:eq:el-cpi} conserva la
suma completa de la cola geométrica
\begin{equation}
 \frac{D_A\alpha^7}{1-D_A\alpha}
 =\sum_{n=0}^{\infty}D_A^{n+1}\alpha^{n+7}.
 \label{ii:eq:el-cola-regional}
\end{equation}
No es una truncación que pueda modificarse silenciosamente al cambiar la
precisión. Tras el término de índice \(N\), la cola restante es
\(D_A^{N+2}\alpha^{N+8}/(1-D_A\alpha)\), lo que permite controlar la
evaluación sin acudir a una magnitud objetivo.

\begin{definition}[Secciones de acción y lector de retorno]
Las secciones anterior y posterior al retorno son
\begin{equation}
 \hbar_{\mathrm{pre}}=H_5\,10^{-34}\mathcal U_S,
 \qquad
 \hbar_{\mathrm{ret}}=\eta_{\mathrm{ret}}\,10^{-34}\mathcal U_S.
 \label{ii:eq:el-secciones-accion}
\end{equation}
Siempre que \(\eta_{\mathrm{ret}}>0\), se definen
\begin{equation}
 \delta_{\mathrm{ret}}=H_5-\eta_{\mathrm{ret}},\qquad
 R_{\mathrm{act}}=\frac{\hbar_{\mathrm{pre}}}{\hbar_{\mathrm{ret}}}
 =\frac{H_5}{\eta_{\mathrm{ret}}}.
 \label{ii:eq:el-ratio}
\end{equation}
\end{definition}

\begin{proposition}[Positividad e invariancia de la razón]
En los intervalos usados en la prueba de \eqref{ii:eq:el-decada},
\(0<\eta_{\mathrm{ret}}<H_5\), y por tanto
\(R_{\mathrm{act}}>1\). La razón no cambia al sustituir la base
\(\mathcal U_S\) por cualquier otra base positiva común.
\end{proposition}
\begin{proof}
Los términos de \(\mathcal C_\pi\) son positivos. Para una cota
elemental basta usar \(\alpha<0.008\), \(\varphi<1.619\),
\(\alpha>0.007\), \(D_A<1\) y \(\pi>3\). La raíz en
\eqref{ii:eq:el-h5} es mayor que uno, mientras la suma de los términos
negativos es menor que \(0.008001\); en particular, \(H_5>0.99\).
Asimismo,
\[
 0<\frac{90}{\pi}\mathcal C_\pi
 <30\left(169(0.008)^6+
           \frac{(0.008)^7}{1-0.008}\right)<10^{-6}.
\]
Así \(\eta_{\mathrm{ret}}>0.989999>0\) y
\(\eta_{\mathrm{ret}}<H_5\). Finalmente, el cambio de base divide
ambas coordenadas de acción por el mismo factor positivo, que cancela en
su cociente.
\end{proof}

La evaluación posterior proporciona
\begin{align*}
 H_5&=1.0545718176461544696\ldots,\\
 \eta_{\mathrm{ret}}&=1.0545718169150390611\ldots,\\
 R_{\mathrm{act}}&=1.0000000006932817630\ldots.
\end{align*}
La diferencia y el cociente son dos lecturas de las mismas secciones,
una aditiva y otra multiplicativa. No representan dos constantes de
Planck independientes. El paso de acción angular a acción por vuelta
completa se expresa, en cada sección, mediante
\(h=2\pi\hbar\). La posterior carta
\(\mathcal U_S\mapsto\mathrm{J\,s}\) asigna una unidad, pero no
retroactúa sobre el orden decimal ya demostrado.

El retorno regional puede escribirse también en una carta angular.
Esta reparametrización no es una dependencia adicional indispensable:
\eqref{ii:eq:el-h5}--\eqref{ii:eq:el-eta} suministran directamente
\(H_5\), \(\eta_{\mathrm{ret}}\) y \(R_{\mathrm{act}}\). Si se
introducen ángulo y contraángulo, ambos deben transportarse en la misma
carta; no se mezclan grados y radianes dentro de una diferencia.


% BEGIN OWNER /Users/ruben/Documents/ChatGPT/jueces y controles/REVISION_ARGUMENTAL_APERTURAS_CIERRES_20260915/II/manuscript_es/sections/revision_planck.tex
\subsubsection{Carácter de acción y procedencia de sus coeficientes}
\label{ii:sec:revision-planck}

La constante de Planck reducida, \(\hbar\), es el objeto físico que
motiva la realización de la sección de acción angular. Su posición en
la dependencia electrónica debe especificarse en dos niveles. La norma
determinantal \(\Sigma_H\) determina la década de la escala. El carácter
\(H_5\) y el retorno regional determinan después las coordenadas de las
secciones de acción. Esta separación permite exponer el resultado
necesario para el electrón sin anticipar los desarrollos completos de
otros sectores de la mecánica dimensional.

La procedencia del carácter de quinto orden puede explicitarse mediante
invariantes ya fijados en el soporte. En la carta central se considera
\[
 B_c=\begin{pmatrix}7&2\\2&7\end{pmatrix},
 \qquad \lambda_+=9,\quad\lambda_-=5.
\]
El calendario tiene longitud \(12\), la órbita del paso \(3\) tiene
longitud \(4\), y el cociente censal pertinente es
\(104976/1944=54\). Los coeficientes del carácter satisfacen
\begin{equation}
 \frac9{16}=\frac{\lambda_+}{4^2},\quad
 \frac59=\frac{\lambda_-}{\lambda_+},\quad
 \frac7{48}=\frac{\operatorname{tr}(B_c)/2}{4\cdot12},\quad
 \frac1{54}=\frac{1944}{104976}.
 \label{ii:eq:revision-h5-coeficientes}
\end{equation}
La asignación de estos invariantes a los grados \(2,3,4,5\), con la
alternancia de signos de \eqref{ii:eq:el-h5}, pertenece a la definición del
carácter analítico utilizado. Explicitarla es esencial: el conjunto de
cardinales, considerado sin el carácter y sin el orden de sus grados,
admitiría otras expresiones. El contenido constructivo reside en la
composición del soporte, la orientación y la regla analítica declarada.

\subsubsection{Continuación regional y ecuación de renovación}

El selector \(169\) determina el impulso de grado seis. La prolongación
posterior se organiza por la transferencia \(D_A\) y una normalización
unitaria de la línea determinantal. Estas reglas admiten una formulación
en series formales antes de sustituir \(z=\alpha\), con \(D_A\) ya
construido. Escribamos
\[
 \mathscr C(z)=169z^6+R(z),\qquad R(z)\in z^7\mathbb R[[z]].
\]
La concatenación de una prolongación aumenta el grado en una unidad y
multiplica su peso por \(D_A\). La primera prolongación estricta tiene
peso \(D_A\). Por tanto, su ecuación de renovación es
\begin{equation}
 R(z)=D_Az^7+D_AzR(z).
 \label{ii:eq:revision-renovacion}
\end{equation}

\begin{proposition}[Unicidad formal de la continuación normalizada]
La ecuación \eqref{ii:eq:revision-renovacion} determina una única serie y
\[
 \mathscr C(z)=169z^6+\frac{D_Az^7}{1-D_Az}
             =169z^6+\sum_{n=7}^{\infty}D_A^{n-6}z^n.
\]
La realización es analítica en \(|D_Az|<1\).
\end{proposition}
\begin{proof}
El elemento \(1-D_Az\) tiene término independiente uno y es invertible
en el anillo de series formales. Resolver la ecuación produce la serie
exhibida. Equivalentemente, el coeficiente de grado siete es \(D_A\)
y cada coeficiente posterior es \(D_A\) veces el anterior. La serie
geométrica converge absolutamente en el disco indicado.
\end{proof}

La normalización de la primera prolongación es un dato operativo
indispensable. Si se conservaran únicamente el impulso de grado seis y
la razón de concatenación, las series
\[
 169z^6+\lambda\frac{D_Az^7}{1-D_Az}
\]
seguirían compartiendo ambos datos para cualquier \(\lambda\).
La regla unitaria fija \(\lambda=1\) antes de efectuar la evaluación.
Así queda localizada la contribución de cada regla a la unicidad del
lector, y la expresión racional conserva la continuación completa.

\subsubsection{Constante de Planck reducida y retorno de la sección de acción}

En la carta dimensional de \eqref{ii:eq:el-secciones-accion}, la sección
anterior al retorno es la construcción del modelo destinada al contraste
con la constante de Planck reducida:
\[
 \hbar_{\mathrm{pre}}=H_5\,10^{-34}\mathcal U_S.
\]
La sección posterior conserva, además, la compensación regional
\(90\mathscr C(\alpha)/\pi\). Ambas pertenecen a la misma recta de
acción y su cociente \(R_{\mathrm{act}}\) mide el efecto multiplicativo
del retorno. El paso a \(h\) mediante \(2\pi\) utiliza la relación
entre la parametrización angular y una vuelta completa.

Esta construcción interviene efectivamente en
\(R_{\mathrm{act}}^{80-54\alpha+6\Delta_4}\). El electrón conserva
su coordenada central y su trayectoria; la escala cronológica no
reemplaza esa estructura por un cuanto indiferenciado. Tampoco se
multiplica de nuevo por \(10^{-34}\) una masa ya expresada en MeV:
la década pertenece a la sección de acción, cancela en su cociente
adimensional y reaparece únicamente donde actúa la realización
dimensional correspondiente.

El contraste numérico de \(H_5\), de la sección retornada y de
\(h/(2\pi)\) se presenta en la sección~\ref{ii:sec:revision-planck-contraste}.
Así se distingue la denominación del objeto físico, la ecuación que el
modelo le asigna y la precisión con la que dicha ecuación lo reproduce.

% END OWNER

